%%%%%%%%%%  scan idx 348  |  folio 335  %%%%%%%%%%
\(f'\),
évidemment unique.

\textbf{Proposition 4.8.} Soit \(n\) un entier \(\ge0\)~; dans le \(\lambda\)-anneau libre engendré sur \(K\) par une indéterminée \(X\), l'idéal engendré par les \(\lambda^k X\), \(k>n\), est un \(\lambda\)-idéal.

Posons \(A=K[X,\lambda^k X]_{k\ge2}\), et soit \(J\) l'idéal engendré par les \(\lambda^k X\) pour \(k>n\)~; posons \(B=A/J\). Il faut montrer que l'homomorphisme composé : \(A\xrightarrow{\ \lambda_t\ }\Ghat(A)\longrightarrow\Ghat(B)\) s'annule sur \(J\), i.e. que \(\lambda_t(\lambda^k X)\), \(k>n\), a pour image \(1\) dans \(\Ghat(B)\). Or \(\lambda_t(\lambda^k X)=\lambda^k(\lambda_t(X))\), et \(\Ghat(A)\to\Ghat(B)\) est un \(\lambda\)-homomorphisme~; il suffit donc de montrer que pour tout anneau \(C\) et toute série \(\varphi\in\Ghat(C)\), dont les coefficients sont nuls en degrés \(>n\), \(\lambda^k(\varphi)=1\) si \(k>n\). Mais il suffit pour prouver cette dernière propriété de la montrer pour \(C=K[X_1,\ldots,X_n]\), et \(\varphi=1+X_1t+\ldots+X_nt^n\). Plongeant \(K[X_1,\ldots,X_n]\) dans \(K[T_1,\ldots,T_n]\) grâce aux fonctions symétriques, on peut écrire : \(1+X_1t+\ldots+X_nt^n=\prod_{i=1}(1+T_it)\), d'où
\[
  \lambda_u(1+X_1t+\ldots+X_nt^n)=\prod_{i=1}^{n}\lambda_u(1+T_it)=\prod_{i=1}^{n}\bigl(\underline{1}+\{1+T_it\}u\bigr),
\]
qui est de degré \(n\) en \(u\).

\textbf{Corollaire 4.9.} Le \(\lambda\)-anneau engendré sur \(K\) par le générateur \(X\) soumis aux relations \(\lambda^k(X)=0\) pour \(k>n\), est isomorphe à \(K[X,\ldots,\lambda^n X]\).

\textbf{4.9.1.} On aurait les mêmes énoncés en remplaçant les \(\lambda\) par les \(\gamma\). On peut également remplacer \(X\) par une famille d'indéterminées \((X_i)_{i\in I}\), soumises aux relations \(\lambda^k(X_i)=0\) pour \(k>n_i\), les \(n_i\) étant des entiers \(\ge2\)~; le \(\lambda\)-anneau obtenu est alors \(K[X_i,\lambda^k X_i]_{i\in I,\ 2\le k\le n_i}\).

\textbf{4.10.} Soient \(K\) un anneau binômial, et \(A\) une \(K\)-\(\lambda\)-algèbre augmentée. On désigne par \(B=A[X,\gamma^k X]_{k\ge2}\) le \(\lambda\)-anneau libre engendré par \(X\) sur \(A\)~; 
